\documentclass[12pt]{article}
\usepackage[utf8]{inputenc}
\usepackage[spanish]{babel}
\usepackage{amsmath}
\usepackage{amsfonts}
\usepackage{amssymb}
\usepackage{geometry}
\usepackage{tabularx}
\usepackage{multirow}
\usepackage{longtable}
\usepackage{booktabs}
\usepackage{color}
\usepackage{listings}
\usepackage{hyperref}
\usepackage{url}

\geometry{top=2.5cm, bottom=2.5cm, left=2.5cm, right=2.5cm}

% Definir color para los comandos
\definecolor{codegreen}{rgb}{0,0.6,0}
\definecolor{codegray}{rgb}{0.5,0.5,0.5}
\definecolor{codepurple}{rgb}{0.58,0,0.82}

% Configuración para código ensamblador
\lstdefinelanguage{MIPS}{
    morekeywords={add,addi,sub,subi,mul,div,beq,bne,j,jr,jal,lw,sw,li,la,addiu,addu},
    morekeywords=[2]{syscall,move,nop},
    sensitive=false,
    morecomment=[l]{\#},
    morestring=[b]"
}

\lstset{
    language=MIPS,
    basicstyle=\ttfamily\small,
    keywordstyle=\color{blue},
    keywordstyle=[2]\color{codepurple},
    commentstyle=\color{codegreen},
    stringstyle=\color{codegray},
    numbers=left,
    numberstyle=\tiny\color{codegray},
    stepnumber=1,
    numbersep=5pt,
    frame=single,
    breaklines=true
}

\title{\textbf{Guía de Referencia}\\MIPS32 - Arquitectura y Conjunto de Instrucciones}
\author{Estudiante de Arquitectura de Computadores}
\date{\today}

\begin{document}

\maketitle

\tableofcontents
\newpage

\section{Introducción a MIPS32}

MIPS (Microprocessor without Interlocked Pipeline Stages) es una arquitectura de tipo RISC (Reduced Instruction Set Computer) desarrollada originalmente por John L. Hennessy en la Universidad de Stanford. El libro "Computer Organization and Design" de Patterson y Hennessy utiliza MIPS como arquitectura de estudio debido a su elegancia y simplicidad \cite{hennessy2012computer}.

\subsection{Características Principales}
\begin{itemize}
    \item Conjunto de instrucciones reducido y ortogonal
    \item Formato de instrucciones de 32 bits de longitud fija
    \item 32 registros de propósito general de 32 bits
    \item Modos de direccionamiento simples
    \item Ideal para la enseñanza de arquitectura de computadores
\end{itemize}

\section{Convención de Notaciones}

En este documento se utilizarán las siguientes notaciones estándar para describir las instrucciones MIPS \cite{mipsguide}:

\begin{center}
    \begin{tabularx}{0.8\textwidth}{|l|X|}
        \hline
        \textbf{Notación} & \textbf{Descripción} \\
        \hline
        Rd, Rs, Rt & Registros (destino, fuente1, fuente2) \\
        \hline
        Imm16 & Valor inmediato de 16 bits \\
        \hline
        Label & Etiqueta que representa una dirección de memoria \\
        \hline
        Offset(Rb) & Direccionamiento base+desplazamiento \\
        \hline
        \# comentario & Comentario en ensamblador \\
        \hline
    \end{tabularx}
\end{center}

\section{Registros del MIPS32}

La arquitectura MIPS32 dispone de 32 registros de propósito general de 32 bits cada uno. La siguiente tabla muestra la convención de nombres y el uso estándar de cada registro \cite{spimref}:

\begin{longtable}{|c|c|c|p{5cm}|}
    \hline
    \textbf{Número} & \textbf{Nombre} & \textbf{Nombre Alternativo} & \textbf{Uso} \\
    \hline
    0 & \$zero & \$0 & Siempre contiene el valor 0 \\
    \hline
    1 & \$at & \$1 & Reservado para el ensamblador \\
    \hline
    2 & \$v0 & \$2 & Valor de retorno de funciones \\
    \hline
    3 & \$v1 & \$3 & Valor de retorno de funciones \\
    \hline
    4 & \$a0 & \$4 & Argumento 1 para funciones \\
    \hline
    5 & \$a1 & \$5 & Argumento 2 para funciones \\
    \hline
    6 & \$a2 & \$6 & Argumento 3 para funciones \\
    \hline
    7 & \$a3 & \$7 & Argumento 4 para funciones \\
    \hline
    8-15 & \$t0-\$t7 & \$8-\$15 & Temporales (no preservados entre llamadas) \\
    \hline
    16-23 & \$s0-\$s7 & \$16-\$23 & Registros preservados (callee-saved) \\
    \hline
    24-25 & \$t8-\$t9 & \$24-\$25 & Temporales adicionales \\
    \hline
    26-27 & \$k0-\$k1 & \$26-\$27 & Reservados para el kernel del SO \\
    \hline
    28 & \$gp & \$28 & Puntero al área global \\
    \hline
    29 & \$sp & \$29 & Puntero de pila (Stack Pointer) \\
    \hline
    30 & \$fp & \$30 & Puntero de frame (Frame Pointer) \\
    \hline
    31 & \$ra & \$31 & Dirección de retorno \\
    \hline
    \caption{Registros de propósito general MIPS32 \cite{regusage}}
    \label{tab:registros}
\end{longtable}

\subsection{Convenciones de Uso de Registros}

\begin{itemize}
    \item \textbf{\$zero}: Siempre lee como 0, las escrituras no tienen efecto
    \item \textbf{\$v0-\$v1}: Contienen el resultado de una función
    \item \textbf{\$a0-\$a3}: Se usan para pasar los primeros 4 argumentos a funciones
    \item \textbf{\$sp}: Apunta al tope de la pila (crece hacia direcciones menores)
    \item \textbf{\$ra}: Contiene la dirección de retorno al hacer \texttt{jal}
\end{itemize}

\section{Conjunto de Instrucciones}

\subsection{Instrucciones Aritméticas}

Las operaciones aritméticas básicas siguen el formato: \texttt{operación Rd, Rs, Rt} donde \texttt{Rd = Rs operación Rt} \cite{mipsguide}.

\begin{lstlisting}[caption=Ejemplos de instrucciones aritméticas, label=lst:aritmeticas]
add $t0, $t1, $t2    # $t0 = $t1 + $t2
sub $t0, $t1, $t2    # $t0 = $t1 - $t2
addi $t0, $t1, 10    # $t0 = $t1 + 10
mul $t0, $t1, $t2    # $t0 = $t1 * $t2
\end{lstlisting}

\begin{longtable}{|l|l|l|}
    \hline
    \textbf{Instrucción} & \textbf{Sintaxis} & \textbf{Descripción} \\
    \hline
    ADD & add Rd, Rs, Rt & Rd = Rs + Rt (con detección de overflow) \\
    \hline
    ADDI & addi Rt, Rs, Imm16 & Rt = Rs + Imm16 (con overflow) \\
    \hline
    ADDU & addu Rd, Rs, Rt & Rd = Rs + Rt (sin overflow) \\
    \hline
    ADDIU & addiu Rt, Rs, Imm16 & Rt = Rs + Imm16 (sin overflow) \\
    \hline
    SUB & sub Rd, Rs, Rt & Rd = Rs - Rt (con overflow) \\
    \hline
    SUBU & subu Rd, Rs, Rt & Rd = Rs - Rt (sin overflow) \\
    \hline
    MUL & mul Rd, Rs, Rt & Rd = Rs * Rt \\
    \hline
    DIV & div Rs, Rt & Lo = Rs / Rt, Hi = Rs \% Rt \\
    \hline
    \caption{Instrucciones aritméticas MIPS32}
\end{longtable}

\subsection{Instrucciones Lógicas}

\begin{lstlisting}[caption=Ejemplos de operaciones lógicas, label=lst:logicas]
and $t0, $t1, $t2     # $t0 = $t1 AND $t2
or  $t0, $t1, $t2     # $t0 = $t1 OR $t2
nor $t0, $t1, $t2     # $t0 = NOT($t1 OR $t2)
xor $t0, $t1, $t2     # $t0 = $t1 XOR $t2
\end{lstlisting}

\begin{longtable}{|l|l|l|}
    \hline
    \textbf{Instrucción} & \textbf{Sintaxis} & \textbf{Descripción} \\
    \hline
    AND & and Rd, Rs, Rt & Rd = Rs \& Rt \\
    \hline
    ANDI & andi Rt, Rs, Imm16 & Rt = Rs \& Imm16 \\
    \hline
    OR & or Rd, Rs, Rt & Rd = Rs | Rt \\
    \hline
    ORI & ori Rt, Rs, Imm16 & Rt = Rs | Imm16 \\
    \hline
    NOR & nor Rd, Rs, Rt & Rd = $\sim$(Rs | Rt) \\
    \hline
    XOR & xor Rd, Rs, Rt & Rd = Rs \^ Rt \\
    \hline
    XORI & xori Rt, Rs, Imm16 & Rt = Rs \^ Imm16 \\
    \hline
    \caption{Instrucciones lógicas MIPS32}
\end{longtable}

\subsection{Instrucciones de Desplazamiento}

\begin{longtable}{|l|l|l|}
    \hline
    \textbf{Instruccion} & \textbf{Sintaxis} & \textbf{Descripcion} \\
    \hline
    SLL & sll Rd, Rt, Imm5 & Rd = Rt $\ll$ Imm5 (desplazamiento logico izquierda) \\
    \hline
    SLLV & sllv Rd, Rt, Rs & Rd = Rt $\ll$ Rs \\
    \hline
    SRL & srl Rd, Rt, Imm5 & Rd = Rt $\gg$ Imm5 (desplazamiento logico derecha, rellena con 0) \\
    \hline
    SRA & sra Rd, Rt, Imm5 & Rd = Rt $\gg$ Imm5 (desplazamiento aritmetico, rellena con MSB) \\
    \hline
    \caption{Instrucciones de desplazamiento}
\end{longtable}

\subsection{Instrucciones de Transferencia de Datos}

Estas instrucciones permiten mover datos entre registros y memoria \cite{mipsguide}:

\begin{lstlisting}[caption=Ejemplos de carga y almacenamiento, label=lst:memoria]
lw $t0, 0($s1)        # Carga palabra en $t0 desde Mem[$s1 + 0]
sw $t0, 4($s1)        # Almacena palabra en Mem[$s1 + 4] desde $t0
lb $t0, 0($s1)        # Carga byte (con extension de signo)
lbu $t0, 0($s1)       # Carga byte (sin extension de signo)
li $t0, 100           # Carga inmediata (pseudoinstruccion)
la $t0, etiqueta      # Carga direccion de etiqueta
\end{lstlisting}

\begin{longtable}{|l|l|l|}
    \hline
    \textbf{Instruccion} & \textbf{Sintaxis} & \textbf{Descripcion} \\
    \hline
    LW & lw Rt, Offset16(Rb) & Rt = Mem[Rb + Offset16] (palabra) \\
    \hline
    SW & sw Rt, Offset16(Rb) & Mem[Rb + Offset16] = Rt (palabra) \\
    \hline
    LB & lb Rt, Offset16(Rb) & Rt = Mem[Rb + Offset16] (byte con signo) \\
    \hline
    LBU & lbu Rt, Offset16(Rb) & Rt = Mem[Rb + Offset16] (byte sin signo) \\
    \hline
    LH & lh Rt, Offset16(Rb) & Rt = Mem[Rb + Offset16] (half-word con signo) \\
    \hline
    LUI & lui Rt, Imm16 & Rt = Imm16 $\ll$ 16 \\
    \hline
    \caption{Instrucciones de carga/almacenamiento}
\end{longtable}

\subsection{Instrucciones de Comparación y Salto}

Las instrucciones de comparación (\texttt{slt} - Set on Less Than) establecen un registro a 1 si la condición se cumple, o a 0 en caso contrario \cite{mipsguide}.

\begin{lstlisting}[caption=Ejemplos de comparación y saltos, label=lst:branch]
slt $t0, $t1, $t2     # $t0 = 1 si $t1 < $t2, sino 0
beq $t1, $t2, label   # Salta a label si $t1 == $t2
bne $t1, $t2, label   # Salta a label si $t1 != $t2
j label               # Salto incondicional (pseudoinstrucción)
jal funcion           # Salta a funcion y guarda retorno en $ra
jr $ra                # Retorna de función
\end{lstlisting}

\begin{longtable}{|l|l|l|}
    \hline
    \textbf{Instrucción} & \textbf{Sintaxis} & \textbf{Descripción} \\
    \hline
    SLT & slt Rd, Rs, Rt & Rd = 1 si Rs < Rt, sino 0 (con signo) \\
    \hline
    SLTU & sltu Rd, Rs, Rt & Rd = 1 si Rs < Rt, sino 0 (sin signo) \\
    \hline
    SLTI & slti Rt, Rs, Imm16 & Rt = 1 si Rs < Imm16, sino 0 \\
    \hline
    BEQ & beq Rs, Rt, Label & Salta si Rs == Rt \\
    \hline
    BNE & bne Rs, Rt, Label & Salta si Rs != Rt \\
    \hline
    BGEZ & bgez Rs, Label & Salta si Rs >= 0 \\
    \hline
    BLEZ & blez Rs, Label & Salta si Rs <= 0 \\
    \hline
    BGTZ & bgtz Rs, Label & Salta si Rs > 0 \\
    \hline
    BLTZ & bltz Rs, Label & Salta si Rs < 0 \\
    \hline
    J & j Label & Salto incondicional a Label \\
    \hline
    JAL & jal Label & Salto y guarda retorno en \$ra \\
    \hline
    JR & jr Rs & Salto a dirección contenida en Rs \\
    \hline
    \caption{Instrucciones de comparación y salto}
\end{longtable}

\subsection{Instrucciones de Control del Sistema - Syscall}

La instrucción \texttt{syscall} permite realizar llamadas al sistema operativo (o al simulador en MARS). El valor en \$v0 determina qué operación realizar \cite{mipsguide}:

\begin{longtable}{|c|l|l|}
    \hline
    \textbf{Código (\$v0)} & \textbf{Operación} & \textbf{Argumentos/Resultados} \\
    \hline
    1 & Imprimir entero & \$a0 = entero a imprimir \\
    \hline
    4 & Imprimir cadena & \$a0 = dirección de cadena terminada en null \\
    \hline
    5 & Leer entero & Resultado devuelto en \$v0 \\
    \hline
    8 & Leer cadena & \$a0 = buffer, \$a1 = longitud máxima \\
    \hline
    10 & Salir del programa & - \\
    \hline
    11 & Imprimir carácter & \$a0 = carácter a imprimir \\
    \hline
    12 & Leer carácter & Resultado devuelto en \$v0 \\
    \hline
    \caption{Syscalls útiles en MARS}
\end{longtable}

\begin{lstlisting}[caption=Ejemplo de programa completo con syscalls, label=lst:syscall]
.data
mensaje: .asciiz "Hola Mundo\n"
numero:  .word 42

.text
main:
    # Imprimir cadena
    li $v0, 4
    la $a0, mensaje
    syscall
    
    # Imprimir entero
    li $v0, 1
    lw $a0, numero
    syscall
    
    # Salir
    li $v0, 10
    syscall
\end{lstlisting}

\section{Estructura de un Programa en Ensamblador MIPS}

Un programa típico en MIPS tiene dos secciones principales \cite{mipsguide}:

\begin{lstlisting}[caption=Estructura básica de un programa MIPS, label=lst:estructura]
.data
    # Variables globales y datos inicializados
    mensaje: .asciiz "Ingrese un numero: "
    resultado: .word 0

.text
    .globl main
main:
    # Código del programa
    
    # Salida del programa
    li $v0, 10
    syscall
\end{lstlisting}

\subsection{Directivas de Datos más Comunes}

\begin{longtable}{|l|l|}
    \hline
    \textbf{Directiva} & \textbf{Descripción} \\
    \hline
    .byte valor & Almacena uno o más bytes \\
    \hline
    .half valor & Almacena half-words (16 bits) \\
    \hline
    .word valor & Almacena palabras (32 bits) \\
    \hline
    .asciiz "texto" & Almacena cadena terminada en null \\
    \hline
    .space N & Reserva N bytes sin inicializar \\
    \hline
    \caption{Directivas de datos en MIPS}
\end{longtable}

\section{Pseudoinstrucciones}

MARS (al igual que otros ensambladores MIPS) soporta pseudoinstrucciones que facilitan la programación. El ensamblador las traduce automáticamente a instrucciones reales \cite{mipsguide}:

\begin{longtable}{|l|l|l|}
    \hline
    \textbf{Pseudoinstrucción} & \textbf{Equivalencia} & \textbf{Descripción} \\
    \hline
    li Rd, Imm & addiu Rd, \$zero, Imm & Load Immediate \\
    \hline
    la Rd, Label & addiu Rd, \$zero, direccion & Load Address \\
    \hline
    move Rd, Rs & addiu Rd, Rs, 0 & Copia registro \\
    \hline
    nop & sll \$zero, \$zero, 0 & No operation \\
    \hline
    b Label & beq \$zero, \$zero, Label & Salto incondicional \\
    \hline
    blt Rs, Rt, Label & slt at, Rs, Rt; bne at, \$zero, Label & Branch if less than \\
    \hline
    neg Rd, Rs & sub Rd, \$zero, Rs & Negación \\
    \hline
    \caption{Pseudoinstrucciones comunes}
\end{longtable}

\section{Modos de Direccionamiento}

MIPS32 soporta varios modos de direccionamiento \cite{mipsguide}:

\begin{itemize}
    \item \textbf{Direccionamiento por registro}: El operando está en un registro (ej. \texttt{add \$t0, \$t1, \$t2})
    \item \textbf{Direccionamiento inmediato}: El operando es una constante (ej. \texttt{addi \$t0, \$t1, 10})
    \item \textbf{Direccionamiento base+desplazamiento}: Acceso a memoria usando registro base y offset (ej. \texttt{lw \$t0, 8(\$s1)})
    \item \textbf{Direccionamiento relativo a PC}: Usado por instrucciones de salto condicional (\texttt{beq}, \texttt{bne})
    \item \textbf{Direccionamiento pseudo directo}: Usado por \texttt{j} y \texttt{jal}
\end{itemize}

\section{Ejemplos Prácticos}

\subsection{Suma de dos números ingresados por el usuario}

\begin{lstlisting}[caption=Programa para sumar dos números, label=lst:ejemplo1]
.data
    prompt1: .asciiz "Ingrese el primer numero: "
    prompt2: .asciiz "Ingrese el segundo numero: "
    result:  .asciiz "La suma es: "

.text
    .globl main
main:
    # Solicitar primer numero
    li $v0, 4
    la $a0, prompt1
    syscall
    
    li $v0, 5
    syscall
    move $t0, $v0
    
    # Solicitar segundo numero
    li $v0, 4
    la $a0, prompt2
    syscall
    
    li $v0, 5
    syscall
    move $t1, $v0
    
    # Sumar
    add $t2, $t0, $t1
    
    # Mostrar resultado
    li $v0, 4
    la $a0, result
    syscall
    
    li $v0, 1
    move $a0, $t2
    syscall
    
    # Salir
    li $v0, 10
    syscall
\end{lstlisting}

\subsection{Bucle for para sumar números del 1 al N}

\begin{lstlisting}[caption=Suma de 1 a N, label=lst:ejemplo2]
.data
    prompt: .asciiz "Ingrese N: "
    result: .asciiz "La suma de 1 a N es: "

.text
    .globl main
main:
    li $v0, 4
    la $a0, prompt
    syscall
    
    li $v0, 5
    syscall
    move $t0, $v0          # $t0 = N
    
    li $t1, 0              # suma = 0
    li $t2, 1              # contador = 1
    
bucle:
    bgt $t2, $t0, fin      # if contador > N: salir
    add $t1, $t1, $t2      # suma += contador
    addi $t2, $t2, 1       # contador++
    j bucle
    
fin:
    li $v0, 4
    la $a0, result
    syscall
    
    li $v0, 1
    move $a0, $t1
    syscall
    
    li $v0, 10
    syscall
\end{lstlisting}

\section{Referencias}

\begin{thebibliography}{9}
    \bibitem{hennessy2012computer}
        Patterson, D. A., \& Hennessy, J. L. (2012). \textit{Computer Organization and Design: The Hardware/Software Interface}. 4th ed. Morgan Kaufmann.
    
    \bibitem{mipsguide}
        UNSW Sydney. (2024). \textit{MIPS Instruction Set Guide}. COMP1521 Course Materials.
    
    \bibitem{spimref}
        SPIM Documentation. \textit{MIPS Registers and Usage Convention}. Wheaton College.
    
    \bibitem{regusage}
        MIPS Technologies. (2011). \textit{MIPS32 74Kc Programming Manual}.
\end{thebibliography}

\end{document}
